\chapter{داده‌ها و روش}
\label{ch:method}

\section{منطقه مطالعه}
پنج «محدوده مطالعاتی» در سوی امارات و تنگه هرمز انتخاب شد: فجیره، خورفکان، شارجه/دبی و
راس‌الخیمه (هر کدام مربعی ۲۰ کیلومتری) و آبراه هرمز (۴۰ کیلومتری). همه داده‌های آموزش،
اعتبارسنجی و آزمون از همین پنج محدوده آمد. برای سنجش تعمیم به منطقه‌ای تازه، هفت محدوده
۲۰ کیلومتری در ساحل ایران نیز تعریف شد: بندرعباس، لارک/هرمز، کیش، عسلویه، بوشهر، خارک و
جاسک (شکل~\ref{fig:map}). مدل هرگز روی داده این هفت محدوده آموزش ندید و پیکربندی‌اش بر پایه
آن‌ها انتخاب نشد.

\begin{figure}[htbp]
\centering
\begin{tikzpicture}
\begin{axis}[
  width=0.95\textwidth, height=0.85\textwidth,
  xmin=47.5, xmax=60, ymin=23, ymax=30.6,
  unit vector ratio=0.895 1,
  unbounded coords=jump,
  axis background/.style={fill=blue!4},
  xlabel={\rl{طول جغرافیایی (درجه شرقی)}},
  ylabel={\rl{عرض جغرافیایی (درجه شمالی)}},
  tick label style={font=\footnotesize},
]
\addplot[gray!60!black, thin] table[x=lon, y=lat, col sep=comma] {../data/coast.csv};
\addplot[red!80!black, thick] table[x=lon, y=lat, col sep=comma] {../data/boxes.csv};
\node[font=\small, gray!40!black] at (axis cs:51.6,26.9) {\rl{خلیج فارس}};
\node[font=\small, gray!40!black] at (axis cs:58.6,24.3) {\rl{دریای عمان}};
\node[font=\scriptsize] at (axis cs:56.25,27.45) {\rl{بندرعباس}};
\node[font=\scriptsize] at (axis cs:50.3,29.65) {\rl{خارک}};
\node[font=\scriptsize] at (axis cs:53.98,26.15) {\rl{کیش}};
\node[font=\scriptsize] at (axis cs:57.8,25.2) {\rl{جاسک}};
\node[font=\scriptsize] at (axis cs:56.9,24.95) {\rl{فجیره}};
\node[font=\scriptsize] at (axis cs:55.3,25.05) {\rl{دبی}};
\end{axis}
\end{tikzpicture}
\caption[منطقه مطالعه]{منطقه مطالعه: پنج محدوده مطالعاتی امارات و تنگه هرمز (داده‌های آموزش و آزمون) و
هفت محدوده ساحل ایران (منطقه نگه‌داشته برای سنجش تعمیم). خط ساحل از \lr{OpenStreetMap}.}
\label{fig:map}
\end{figure}

\section{داده‌ها}
\begin{itemize}
\item \textbf{تصاویر راداری:} محصولات \lr{IW GRDH} ماهواره‌های \lr{Sentinel-1A}
(سال‌های ۲۰۲۴ و ۲۰۲۵ و نیمه نخست ۲۰۲۶) و \lr{Sentinel-1C/1D} (سپتامبر ۲۰۲۶) از \lr{CDSE}
\cite{cdse}؛ از هر تصویر تنها پنجره پیرامون محدوده خوانده شد.
\item \textbf{آشکارسازی‌های راداری و حضور \lr{AIS}:} از \lr{API} عمومی \lr{GFW} \cite{gfwapi}، در
خانه‌های \LRE{0٫01} درجه‌ای و ساعتی.
\item \textbf{خشکی:} چندضلعی‌های خشکی \lr{OpenStreetMap} \cite{osm}، با نوار حائل ۱۵۰ متری در دریا.
\item \textbf{سکوهای دریایی:} مجموعه \lr{Offshore Platform Dataset} \cite{spanier2026}، با طول عمر
هر سکو، تا سکوهای موجود در تاریخ هر تصویر پوشانده شوند.
\end{itemize}

\section{برچسب‌گذاری خودکار}
\label{sec:labels}
برچسب‌گذاری دستی هزاران شناور ممکن نبود؛ برچسب‌ها از کنار هم گذاشتن آشکارسازی‌های
\lr{CFAR} ما با آشکارسازی‌های \lr{GFW} در همان تصویر ساخته شد. چون \lr{GFW} داده را به‌صورت
شمارش در خانه‌های حدوداً یک‌کیلومتری منتشر می‌کند، آشکارسازی ما تنها وقتی برچسب گرفت که
خانه‌اش بی‌ابهام بود: دقیقاً یک شناور تطبیق‌یافته با \lr{AIS}، هیچ آشکارسازی بدون \lr{AIS}
و هیچ آشکارسازی دیگری از ما. آشکارسازی‌های نزدیک لبه خانه، اگر خانه همسایه چیزی داشت،
مبهم شمرده شدند. از هر آشکارسازی برشی ۱۲۸ پیکسلی (دو قطبش) ساخته شد. افزون بر مثبت‌ها،
برش‌های منفی از دریای خالی و برش‌هایی با ناحیه «نادیده» (اجسامی که برچسب مطمئن نداشتند)
ساخته شد؛ آشکارسازی بی‌برچسب هرگز منفی شمرده نشد. تقسیم آموزش، اعتبارسنجی و آزمون بر پایه
تصویر انجام شد، نه تصادفی، تا برش‌های مجاورِ یک تصویر میان بخش‌ها نشت نکنند
(جدول~\ref{tab:data}).

\begin{table}[htbp]
\centering
\small
\begin{adjustbox}{max width=\textwidth}
\begin{tabular}{lcccccc}
\toprule
بخش & تصویر & گذر-محدوده & مثبت با \lr{AIS} & مثبت بدون \lr{AIS} & منفی & نادیده \\
\midrule
آموزش & \LRE{329} & \LRE{649} & \LRE{17٬701} & \LRE{1٬603} & \LRE{53٬465} & \LRE{14٬293} \\
اعتبارسنجی & \LRE{51} & \LRE{103} & \LRE{2٬893} & \LRE{224} & \LRE{8٬302} & \LRE{2٬200} \\
آزمون & \LRE{42} & \LRE{92} & \LRE{2٬679} & \LRE{144} & \LRE{7٬307} & \LRE{2٬106} \\
\bottomrule
\end{tabular}
\end{adjustbox}
\caption[داده‌های آموزش، اعتبارسنجی و آزمون]{داده‌های آموزش، اعتبارسنجی و آزمون در پنج محدوده مطالعاتی (\lr{Sentinel-1A}، ۲۰۲۴ تا ۲۰۲۵). تقسیم بر پایه تصویر است، نه تصادفی. «مثبت بدون \lr{AIS}» آشکارسازی‌های \lr{GFW} بدون تطبیق \lr{AIS} است که در مدل نهایی مثبت گرفته شد.}
\label{tab:data}
\end{table}


\section{مدل}
مدل یک آشکارساز نقطه‌ای چگال است: رمزگذار \lr{EfficientNet-B0} \cite{tan2019} با وزن‌های
از پیش آموزش‌دیده \lr{ImageNet}، رمزگشای \lr{U-Net} \cite{ronneberger2015} و خروجی با گام ۲
(نیمی از تفکیک ورودی). دو سر خروجی دارد: نقشه احتمال مرکز شناور و جابه‌جایی زیرپیکسلی مرکز.
نقشه مرکز با زیان کانونی کاهش‌یافته \cite{lin2017,zhou2019} و جابه‌جایی با زیان \lr{L1} روی
پیکسل‌های مرکز آموزش دید؛ ناحیه‌های «نادیده» در زیان وزن صفر داشتند. سرهای طول، پهنا،
جهت و رده که در طرح اولیه بودند کنار گذاشته شدند، چون برچسب قابل‌اعتمادی برایشان نبود.

آموزش در \lr{Kaggle} روی یک \lr{GPU} مدل \lr{Tesla T4} انجام شد: ۱۵ دوره، هر دوره
\LRE{40٬000} نمونه وزن‌دار، دسته ۶۴ و نرخ یادگیری یک‌چرخه‌ای با بیشینه \LRE{0٫001}.
مدل نخست (\lr{run 1}) آشکارسازی‌های بدون \lr{AIS} را نادیده می‌گرفت. مدل نهایی
(\lr{run 2}) سه تفاوت داشت: آشکارسازی‌های بدون \lr{AIS} \lr{GFW} مثبت شمرده شدند، برش‌های
تصادفی ۹۶ پیکسلی از برش‌های ۱۲۸ پیکسلی به کار رفت تا شناور همیشه در مرکز نباشد، و وزن‌های
همه دوره‌ها نگه داشته شد. مدل به \lr{ONNX} با دقت \lr{fp32} صادر شد؛ کوانتش \lr{INT8}
خروجی‌ها را خراب می‌کرد و کنار گذاشته شد.

\section{استنتاج و انتخاب پیکربندی}
پنجره محدوده به کاشی‌های ۲۵۶ پیکسلی با هم‌پوشانی ۶۴ پیکسل تقسیم می‌شود، نقشه‌ها به هم
دوخته می‌شوند و بیشینه‌های محلی بالاتر از آستانه ۰٫۶۰ آشکارسازی‌اند. دوره آموزش، اندازه
کاشی و آستانه تنها روی ۲۵ گذر اعتبارسنجی و در برابر \lr{GFW} برگزیده شد
(جدول~\ref{tab:validation} در فصل نتایج)، با قاعده‌ای که پیش از دیدن هر عدد آزمون توافق شد:
بیشینه میانگین \lr{F1} دو قاعده تطبیق.

\section{روش سنجش}
\label{sec:evaluation}
برچسب‌های آموزشی خود از \lr{CFAR} آمده‌اند، پس \lr{CFAR} را نمی‌توان در برابر آن‌ها منصفانه
سنجید. مرجع مستقل، همه آشکارسازی‌های راداری \lr{GFW} (با و بدون \lr{AIS}) در محدوده و در ساعت
گذر است. چون \lr{GFW} شمارش هر خانه \LRE{0٫01} درجه‌ای را منتشر می‌کند، آشکارسازی‌ها در سطح
خانه تطبیق می‌یابند: آشکارسازی ما درست است اگر خانه‌اش هنوز آشکارسازی تطبیق‌نیافته‌ای از
\lr{GFW} داشته باشد («همان خانه»)؛ قاعده آسان‌گیرتر، خانه‌های همسایه را نیز می‌پذیرد. دقت،
بازیابی و \lr{F1} از این شمارش‌ها، و بازیابی جداگانه روی آشکارسازی‌های با و بدون \lr{AIS}
محاسبه می‌شود. همه آشکارسازها روی یک پنجره و با یک نقاب خشکی و سکو اجرا می‌شوند. تفاضل
مدل و \lr{CFAR} با بوت‌استرپ جفتی روی گذرها \cite{efron1993} و بازه ۹۵٪ گزارش می‌شود.

دو لایه افزوده، پیش از دیدن نتایج ساحل ایران تعریف شد: امتیازدهی تنها روی خانه‌هایی که مرکزشان
دست‌کم ۱ کیلومتر از خشکی فاصله دارد، زیرا \lr{GFW} در این نوار آشکارسازی گزارش نمی‌کند؛ و
زیرمجموعه گذرهایی که هیچ قابِ گذرشان در آموزش یا اعتبارسنجی نبوده است. \lr{GFW} حقیقت
زمینی نیست: آشکارسازی‌های خودش را دارد و روی \lr{Sentinel-1C/1D} کم‌پشت‌تر است؛ پس اعداد این
پژوهش «سازگاری با یک سامانه منتشرشده مستقل»اند، نه دقت مطلق.

\section{تلفیق با \lr{AIS} و نامزدهای شناور تاریک}
\label{sec:fusion}
تنها \lr{AIS} رایگانِ در دسترس در زمان گذرها، حضور ساعتی \lr{GFW} در خانه‌های \LRE{0٫01}
درجه‌ای بود (مسیر نقطه‌به‌نقطه شناورها در دسترس نبود). برای هر گذر، حضور ساعت پیش، ساعت گذر
و ساعت پس از آن خوانده می‌شود و برای هر شناور «دروازه‌ای» ساخته می‌شود: خانه‌های حضورش،
گسترده به اندازه نیم‌قطر خانه، به‌علاوه جابه‌جایی ظاهری شناور متحرک در تصویر رادار
(حاصل‌ضرب کران \lr{R/V} برابر ۱۱۸ ثانیه در سرعت)، به‌علاوه یک حاشیه. آشکارسازی‌ها و شناورها
با تخصیص یک‌به‌یک \cite{kuhn1955} جفت می‌شوند و آشکارسازی‌ای که در هیچ دروازه‌ای نیفتد
«تاریک» است. پیش از آن، آشکارسازی‌های تکراری (کمتر از ۱۰۰ متر از هم)، نزدیک‌تر از ۱ کیلومتر
به خشکی، نزدیک‌تر از ۳۰۰ متر به سکو و نزدیک‌تر از ۵۰۰ متر به لبه محدوده کنار گذاشته
می‌شوند. حاشیه صفر متر، روی گذرهای اعتبارسنجی و با بیشترین سازگاری با تطبیق خود \lr{GFW}
برگزیده شد؛ سطح اطمینان هر نامزد از این است که تا چه حاشیه‌ای همچنان تاریک می‌ماند (زیاد:
حتی با ۲ کیلومتر). گذرهایی که حضور \lr{AIS} در آن‌ها آشکارا ناقص است از نرخ تاریک کنار
گذاشته می‌شوند. هر نامزد همراه «کارت شواهد» ذخیره می‌شود: فاصله تا نزدیک‌ترین دروازه،
داوری \lr{GFW} درباره همان خانه، فاصله تا ساحل و نزدیک‌ترین شناورهای \lr{AIS}.

\section{نقطه‌های پایدار}
\label{sec:persistence}
سازه‌های ثابت (اسکله، بویه، مهاربند، سکوی ثبت‌نشده) و شناورهایی که ماه‌ها لنگر انداخته‌اند
در گذرهای پیاپی در یک نقطه دیده می‌شوند. «نقطه پایدار» با قاعده‌ای که پیش از دیدن نتایج
تعیین شد تعریف شد: جایی که مدل در شعاع ۱۵۰ متر، دست‌کم در ۳ گذر و دست‌کم در ۳۰٪ گذرهای
بایگانی آن محدوده چیزی آشکار کرده است. بایگانی، گذرهای \lr{Sentinel-1A} پیش از ژوئیه ۲۰۲۶ بود
و اثر آن روی گذرهای سپتامبر ۲۰۲۶ سنجیده شد. کنسول این آشکارسازی‌ها را تنها علامت می‌زند و
پنهان نمی‌کند، چون شناورِ لنگرانداخته بدون \lr{AIS} خود برای تحلیلگر مهم است.

\section{سامانه و کنسول}
زنجیره کامل بدون دخالت دستی کار می‌کند: جست‌وجوی تصویر در \lr{CDSE}، خواندن پنجره،
آشکارسازی با مدل، برش تصویر راداری هر آشکارسازی (شمال بالا، ۲٫۵۶ کیلومتر)، ذخیره در پایگاه
داده، تلفیق با \lr{AIS} چند روز پس از گذر، و نمایش در کنسول وب. کنسول فارسی و راست‌به‌چپ
است، تاریخ‌ها را به تقویم خورشیدی نشان می‌دهد و هر آشکارسازی را با تصویر راداری‌اش، سن گذر،
امتیاز مدل و اینکه دقت مدل در آن نقطه سنجیده شده یا نه نشان می‌دهد. یک خط زمان، گذرهای
هر محدوده را نشان می‌دهد و می‌توان نقشه را «در لحظه‌ای از گذشته» دید.
